\documentclass[10pt,letterpaper]{article}
\usepackage[utf8]{inputenc}
\usepackage[intlimits]{amsmath}
\usepackage{amsfonts}
\usepackage{amssymb}
\usepackage{ragged2e}
\usepackage[letterpaper, margin=1in]{geometry}
\usepackage{graphicx}
\usepackage{cancel}
\usepackage{mathtools}
\usepackage{tabularx}
\usepackage{arydshln}
\usepackage{tensor}
\usepackage{array}
\usepackage{xcolor}
\usepackage[boxed]{algorithm}
\usepackage[noend]{algpseudocode}
\usepackage{listings}
\usepackage{textcomp}
% \usepackage[pdf,tmpdir,singlefile]{graphviz}
\usepackage{mathrsfs}
\usepackage{bbm}
\usepackage{tikz}
\usepackage{tikz-cd}
\usepackage{enumitem}
\usepackage{arydshln}
\usepackage{relsize}
\usepackage{multicol}
\usepackage{scalerel}

\usetikzlibrary{bayesnet}
\setlist{noitemsep}

%%%%%%%%%%%%%%%%%%%%%%%%%%%%%
% Formatting commands
%%%%%%%%%%%%%%%%%%%%%%%%%%%%%
\newcommand{\n}{\hfill\break}
\newcommand{\up}{\vspace{-\baselineskip}}
\newcommand{\hangblock}[2]{\par\noindent\settowidth{\hangindent}{\textbf{#1: }}\textbf{#1: }\!\!\!#2}
\newcommand{\lemma}[1]{\hangblock{Lemma}{#1}}
\newcommand{\defn}[1]{\hangblock{Defn}{#1}}
\newcommand{\thm}[1]{\hangblock{Thm}{#1}}
\newcommand{\cor}[1]{\hangblock{Cor}{#1}}
\newcommand{\prop}[1]{\hangblock{Prop}{#1}}
\newcommand{\ex}[1]{\hangblock{Ex}{#1}}
\newcommand{\exer}[1]{\hangblock{Exer}{#1}}
\newcommand{\fact}[1]{\hangblock{Fact}{#1}}
\newcommand{\remark}[1]{\hangblock{Remark}{#1}}
\newcommand{\proven}{\;$\square$\n}
\newcommand{\problem}[1]{\par\noindent{#1}\n}
\newcommand{\problempart}[2]{\par\noindent\indent{}\settowidth{\hangindent}{\textbf{(#1)} \indent{}}\textbf{(#1)} #2\n}
\newcommand{\ptxt}[1]{\textrm{\textnormal{#1}}}
\newcommand{\inlineeq}[1]{\centerline{$\displaystyle #1$}}
\newcommand{\pageline}{\noindent\rule{\textwidth}{0.1pt}}

%%%%%%%%%%%%%%%%%%%%%%%%%%%%%
% Math commands
%%%%%%%%%%%%%%%%%%%%%%%%%%%%%
% Set Theory
\newcommand{\card}[1]{\left|#1\right|}
\newcommand{\set}[1]{\left\{#1\right\}}
\newcommand{\setmid}{\;\middle|\;}
\newcommand{\ps}[1]{\mathcal{P}\left(#1\right)}
\newcommand{\pfinite}[1]{\mathcal{P}^{\ptxt{finite}}\left(#1\right)}
\newcommand{\naturals}{\mathbb{N}}
\newcommand{\N}{\naturals}
\newcommand{\integers}{\mathbb{Z}}
\newcommand{\Z}{\integers}
\newcommand{\rationals}{\mathbb{Q}}
\newcommand{\Q}{\rationals}
\newcommand{\reals}{\mathbb{R}}
\newcommand{\R}{\reals}
\newcommand{\complex}{\mathbb{C}}
\newcommand{\C}{\complex}
\newcommand{\halfPlane}{\mathbb{H}}
\let\H\relax
\newcommand{\H}{\halfPlane}
\newcommand{\comp}{^{\complement}}
\DeclareMathOperator{\Hom}{Hom}
\newcommand{\Ind}{\mathbbm{1}}
\newcommand{\cut}{\setminus}
\DeclareMathOperator{\elem}{elem}

% Graph Theory
\let\deg\relax
\DeclareMathOperator{\deg}{deg}
\newcommand{\degp}{\ptxt{deg}^{+}}
\newcommand{\degn}{\ptxt{deg}^{-}}
\newcommand{\precdot}{\mathrel{\ooalign{$\prec$\cr\hidewidth\hbox{$\cdot\mkern0.5mu$}\cr}}}
\newcommand{\succdot}{\mathrel{\ooalign{$\cdot\mkern0.5mu$\cr\hidewidth\hbox{$\succ$}\cr\phantom{$\succ$}}}}
\DeclareMathOperator{\cl}{cl}
\DeclareMathOperator{\affdim}{affdim}

% Probability
\newcommand{\parSymbol}{\P}
\newcommand{\Prob}{\mathbb{P}}
\renewcommand{\P}{\Prob}
\newcommand{\Avg}{\mathbb{E}}
\newcommand{\E}{\Avg}
\DeclareMathOperator{\Var}{Var}
\DeclareMathOperator{\cov}{cov}
\DeclareMathOperator{\Unif}{Unif}
\DeclareMathOperator{\Binom}{Binom}
\newcommand{\CI}{\mathrel{\text{\scalebox{1.07}{$\perp\mkern-10mu\perp$}}}}

% Standard Math
\newcommand{\inv}{^{-1}}
\newcommand{\abs}[1]{\left|#1\right|}
\newcommand{\ceil}[1]{\left\lceil{}#1\right\rceil{}}
\newcommand{\floor}[1]{\left\lfloor{}#1\right\rfloor{}}
\newcommand{\conj}[1]{\overline{#1}}
\newcommand{\of}{\circ}
\newcommand{\tri}{\triangle}
\newcommand{\inj}{\hookrightarrow}
\newcommand{\surj}{\twoheadrightarrow}
\newcommand{\ndiv}{\nmid}
\renewcommand{\epsilon}{\varepsilon}
\newcommand{\divides}{\mid}
\newcommand{\ndivides}{\nmid}
\DeclareMathOperator{\lcm}{lcm}
\DeclareMathOperator{\sgn}{sgn}
\newcommand{\map}[4]{\!\!\!\begin{array}[t]{rcl}#1 & \!\!\!\!\to & \!\!\!\!#2\\ #3 & \!\!\!\!\mapsto & \!\!\!\!#4\end{array}}
\newcommand{\bigsum}[2]{\smashoperator[lr]{\sum_{\scalebox{#1}{$#2$}}}}

% Linear Algebra
\newcommand{\Id}{\textrm{\textnormal{Id}}}
\newcommand{\im}{\textrm{\textnormal{im}}}
\newcommand{\norm}[1]{\abs{\abs{#1}}}
\newcommand{\tpose}{^{T}\!}
\newcommand{\iprod}[1]{\left<#1\right>}
\DeclareMathOperator{\tr}{tr}
\DeclareMathOperator{\trace}{tr}
\newcommand{\chgBasMat}[3]{\!\!\tensor*[_{#1}]{\left[#2\right]}{_{#3}}}
\newcommand{\vecBas}[2]{\tensor*[]{\left[#1\right]}{_{#2}}}
\DeclareMathOperator{\GL}{GL}
\DeclareMathOperator{\Mat}{Mat}
\DeclareMathOperator{\vspan}{span}
\DeclareMathOperator{\rank}{rank}
\newcommand{\V}[1]{\vec{#1}}
\DeclareMathOperator{\proj}{proj}
\DeclareMathOperator{\compProj}{comp}
\DeclareMathOperator{\row}{row}
\newcommand{\smallPMatrix}[1]{\paren{\begin{smallmatrix}#1\end{smallmatrix}}}
\newcommand{\smallBMatrix}[1]{\brack{\begin{smallmatrix}#1\end{smallmatrix}}}
\newcommand{\pmat}[1]{\begin{pmatrix}#1\end{pmatrix}}
\newcommand{\bmat}[1]{\begin{bmatrix}#1\end{bmatrix}}

% Multilinear Algebra
\newcommand{\Lsym}{\L}
\let\L\relax
\DeclareMathOperator{\L}{\mathscr{L}}
\DeclareMathOperator{\A}{\mathcal{A}}
\DeclareMathOperator{\Alt}{Alt}
\DeclareMathOperator{\Sym}{Sym}
\newcommand{\ot}{\otimes}
\newcommand{\ox}{\otimes}
\DeclareMathOperator{\asc}{asc}
\DeclareMathOperator{\asSet}{set}
\DeclareMathOperator{\sort}{sort}
\DeclareMathOperator{\ringA}{\mathring{A}}

% Topology
\newcommand{\closure}[1]{\overline{#1}}
\newcommand{\uball}{\mathcal{U}}
\DeclareMathOperator{\Int}{Int}
\DeclareMathOperator{\Ext}{Ext}
\DeclareMathOperator{\Bd}{Bd}
\DeclareMathOperator{\rInt}{rInt}
\DeclareMathOperator{\ch}{ch}
\DeclareMathOperator{\ah}{ah}
\newcommand{\LargerTau}{\mathlarger{\mathlarger{\mathlarger{\mathlarger{\tau}}}}}
\newcommand{\Tau}{\mathcal{T}}

% Analysis
\DeclareMathOperator{\Graph}{Graph}
\DeclareMathOperator{\epi}{epi}
\DeclareMathOperator{\hypo}{hypo}
\DeclareMathOperator{\supp}{supp}
\newcommand{\lint}[2]{\underset{#1}{\overset{#2}{{\color{black}\underline{{\color{white}\overline{{\color{black}\int}}\color{black}}}}}}}
\newcommand{\uint}[2]{\underset{#1}{\overset{#2}{{\color{white}\underline{{\color{black}\overline{{\color{black}\int}}\color{black}}}}}}}
\newcommand{\alignint}[2]{\underset{#1}{\overset{#2}{{\color{white}\underline{{\color{white}\overline{{\color{black}\int}}\color{black}}}}}}}
\newcommand{\extint}{\ptxt{ext}\int}
\newcommand{\extalignint}[2]{\ptxt{ext}\alignint{#1}{#2}}
\newcommand{\conv}{\ast}
\newcommand{\pd}[2]{\frac{\partial{}#1}{\partial{}#2}}
\newcommand{\del}{\nabla}
\DeclareMathOperator{\grad}{grad}
\DeclareMathOperator{\curl}{curl}
\let\div\relax
\DeclareMathOperator{\div}{div}
\DeclareMathOperator{\vol}{vol}

% Complex Analysis
\let\Re\relax
\DeclareMathOperator{\Re}{Re}
\let\Im\relax
\DeclareMathOperator{\Im}{Im}
\DeclareMathOperator{\Res}{Res}

% Abstract Algebra
\DeclareMathOperator{\ord}{ord}
\newcommand{\generated}[1]{\left<#1\right>}
\newcommand{\cycle}[1]{\smallPMatrix{#1}}
\newcommand{\id}{\ptxt{id}}
\newcommand{\iso}{\cong}
\DeclareMathOperator{\Aut}{Aut}
\DeclareMathOperator{\SL}{SL}
\DeclareMathOperator{\op}{op}
\newcommand{\isom}[4]{\!\!\!\begin{array}[t]{rcl}#1 & \!\!\!\!\overset{\sim}{\to} & \!\!\!\!#2\\ #3 & \!\!\!\!\mapsto & \!\!\!\!#4\end{array}}
\newcommand{\F}{\mathbb{F}}
\newcommand{\acts}{\;\reflectbox{\rotatebox[origin=c]{-90}{$\circlearrowright$}}\;}
\newcommand{\disjunion}{\mathrel{\text{\scalebox{1.07}{$\perp\mkern-10mu\perp$}}}}
\DeclareMathOperator{\SO}{SO}
\DeclareMathOperator{\stab}{stab}
\DeclareMathOperator{\nullity}{nullity}
\DeclareMathOperator{\Perm}{Perm}
\DeclareMathOperator{\nsubgp}{\vartriangleleft}
\DeclareMathOperator{\notnsubgp}{\ntriangleleft}

% Convex Optimization
\newcommand{\sectionSymbol}{\S}
\let\S\relax
\newcommand{\S}{\mathbb{S}}
\DeclareMathOperator{\dist}{dist}
\DeclareMathOperator{\dom}{dom}
\DeclareMathOperator{\diag}{diag}
\DeclareMathOperator{\ones}{\mathbbm{1}}
\newcommand{\minimizeOver}[3]{\begin{array}{rl}\underset{#1}{\ptxt{minimize}} & #2\\ \ptxt{subject to} & #3\end{array}}
\newcommand{\maximizeOver}[3]{\begin{array}{rl}\underset{#1}{\ptxt{maximize}} & #2\\ \ptxt{subject to} & #3\end{array}}
\newcommand{\minimizationProblem}[2]{\minimizeOver{}{#1}{#2}}
\newcommand{\maximizationProblem}[2]{\maximizeOver{}{#1}{#2}}

% Proofs
\newcommand{\st}{s.t.}
\newcommand{\unique}{!}
\newcommand{\iffdef}{\overset{\ptxt{def}}{\Leftrightarrow}}
\newcommand{\eqdef}{\overset{\ptxt{def}}{=}}
\newcommand{\eqVertical}{\rotatebox[origin=c]{90}{=}}
\newcommand{\mapsfrom}{\mathrel{\reflectbox{\ensuremath{\mapsto}}}}
\newcommand{\mapsdown}{\rotatebox[origin=c]{-90}{$\mapsto$}\mkern2mu}
\newcommand{\mapsup}{\rotatebox[origin=c]{90}{$\mapsto$}\mkern2mu}
\newcommand{\from}{\!\mathrel{\reflectbox{\ensuremath{\to}}}}

% Brackets
\newcommand{\paren}[1]{\left(#1\right)}
\renewcommand{\brack}[1]{\left[#1\right]}
\renewcommand{\brace}[1]{\left\{#1\right\}}
\newcommand{\ang}[1]{\left<#1\right>}

% Algorithms
\algrenewcommand{\algorithmiccomment}[1]{\hskip 1em \texttt{// #1}}
\algrenewcommand\algorithmicrequire{\textbf{Input:}}
\algrenewcommand\algorithmicensure{\textbf{Output:}}
\newcommand{\algP}{\ptxt{\textbf{P}}}
\newcommand{\algNP}{\ptxt{\textbf{NP}}}
\newcommand{\algNPC}{\ptxt{\textbf{NP-Complete}}}
\newcommand{\algNPH}{\ptxt{\textbf{NP-Hard}}}
\newcommand{\algEXP}{\ptxt{\textbf{EXP}}}

%%%%%%%%%%%%%%%%%%%%%%%%%%%%%
% Other commands
%%%%%%%%%%%%%%%%%%%%%%%%%%%%%
\newcommand{\flag}[1]{\textbf{\textcolor{red}{#1}}}
\newcommand{\uSym}{\u}
\let\u\relax
\newcommand{\u}[1]{\underline{#1}}
\newcommand{\bSym}{\b}
\let\b\relax
\newcommand{\b}[1]{\textbf{#1}}
\newcommand{\iSym}{\i}
\let\i\relax
\newcommand{\i}[1]{\textit{#1}}

%%%%%%%%%%%%%%%%%%%%%%%%%%%%%%%%%%%%%%%
% Make l's curvy in math environments %
%%%%%%%%%%%%%%%%%%%%%%%%%%%%%%%%%%%%%%%
\mathcode`l="8000
\begingroup
\makeatletter
\lccode`\~=`\l
\DeclareMathSymbol{\lsb@l}{\mathalpha}{letters}{`l}
\lowercase{\gdef~{\ifnum\the\mathgroup=\m@ne \ell \else \lsb@l \fi}}%
\endgroup

%%%%%%%%%%%%%%%%%%%%%%%%%
% Fix \vdots and \ddots %
%%%%%%%%%%%%%%%%%%%%%%%%%
\usepackage{letltxmacro}
\LetLtxMacro\orgvdots\vdots
\LetLtxMacro\orgddots\ddots

\makeatletter
\DeclareRobustCommand\vdots{%
	\mathpalette\@vdots{}%
}
\newcommand*{\@vdots}[2]{%
	% #1: math style
	% #2: unused
	\sbox0{$#1\cdotp\cdotp\cdotp\m@th$}%
	\sbox2{$#1.\m@th$}%
	\vbox{%
		\dimen@=\wd0 %
		\advance\dimen@ -3\ht2 %
		\kern.5\dimen@
		% remove side bearings
		\dimen@=\wd2 %
		\advance\dimen@ -\ht2 %
		\dimen2=\wd0 %
		\advance\dimen2 -\dimen@
		\vbox to \dimen2{%
			\offinterlineskip
			\copy2 \vfill\copy2 \vfill\copy2 %
		}%
	}%
}
\DeclareRobustCommand\ddots{%
	\mathinner{%
		\mathpalette\@ddots{}%
		\mkern\thinmuskip
	}%
}
\newcommand*{\@ddots}[2]{%
	% #1: math style
	% #2: unused
	\sbox0{$#1\cdotp\cdotp\cdotp\m@th$}%
	\sbox2{$#1.\m@th$}%
	\vbox{%
		\dimen@=\wd0 %
		\advance\dimen@ -3\ht2 %
		\kern.5\dimen@
		% remove side bearings
		\dimen@=\wd2 %
		\advance\dimen@ -\ht2 %
		\dimen2=\wd0 %
		\advance\dimen2 -\dimen@
		\vbox to \dimen2{%
			\offinterlineskip
			\hbox{$#1\mathpunct{.}\m@th$}%
			\vfill
			\hbox{$#1\mathpunct{\kern\wd2}\mathpunct{.}\m@th$}%
			\vfill
			\hbox{$#1\mathpunct{\kern\wd2}\mathpunct{\kern\wd2}\mathpunct{.}\m@th$}%
		}%
	}%
}
\makeatother

\newcommand{\B}{
	\begin{tikzpicture}
	\filldraw [fill=red, draw=black] (0, 0) rectangle (0.37, 0.45);
	\draw [line width=0.5mm, white ] (0.1,0.08) -- (0.1,0.38);
	\draw[line width=0.5mm, white ] (0.1, 0.35) .. controls (0.2, 0.35) and (0.4, 0.2625) .. (0.1, 0.225);
	\draw[line width=0.5mm, white ] (0.1, 0.225) .. controls (0.2, 0.225) and (0.4, 0.1625) .. (0.1, 0.1);
	\end{tikzpicture}
}

\author{Professor Andrew Snowden\\ \small\i{Transcribed by Thomas Cohn}}
\title{Math 493 Lecture 16}
\date{10/30/19} % Can also use \today

\begin{document}
\maketitle
\setlength\RaggedRightParindent{\parindent}
\RaggedRight

\section*{Sylow Theorems}

\par\noindent
Let $G$ be a finite group. Let $p$ be a prime, and write $\card{G}=p^{e}m$, where $p\ndiv{}m$.\n

\defn{
	A \textbf{$p$-Sylow subgroup} of $G$ is a subgroup of order $p^{e}$.\n
}

\thm{
	(First Sylow Theorem) A Sylow subgroup exists, $\forall{}G$, $\forall{}p$.\n
}

\cor{
	(Cauchy's Theorem) If $p\mid\card{G}$, then $G$ has an element of order $p$.\n
	Proof: Let $H$ be a $p$-Sylow subgroup of $H$. $\ord{H}=p^{e}>1$, so let $h\in{}H$ with $h\ne{}1$. Then $\ord(h)\mid\card{H}=p^{e}$, and $\ord(h)\ne{}1$, so $\ord(h)=p^{k}$ for some $k>0$. Thus, $\ord(h^{p^{k-1}})=p$.\proven
}

\par\noindent
Observe: Any conjugate of a Sylow subgroup is a Sylow subgroup.\n

\thm{
	(Second Sylow Theorem) Let $G$ be a group, $H$ a $p$-Sylow subgroup, and let $K$ be any subgroup of $G$. Then $\exists{}H'$ a conjugate of $H$ such that $H'\cap{}K$ is a $p$-Sylow subgroup of $K$.\n
}

\cor{
	Any two $p$-Sylow subgroups in $G$ are conjugate.\n
	Proof: Let $H,K$ be $p$-Sylow subgroups. By the second theorem, there is a conjugate $H'$ of $H$ \st{} $K\cap{}H'$ is a $p$-Sylow subgroup of $K$. Thus, $K\cap{}H'=K$, so $K\subseteq{}H'$. Thus, $K=H'$ (because they're the same order).\proven
}

\cor{
	Any subgroup of $G$ that's a $p$-group is contained in some $p$-Sylow subgroup.\n
	Proof: Let $H$ be a $p$-Sylow subgroup. Let $K$ be a $p$-subgroup. Then there exists a conjugate $H'$ of $H$ \st{} $K\cap{}H'=K$, so $K\subseteq{}H'$ is a $p$-Sylow subgroup.\proven
}

\thm{
	(Third Sylow Theorem) Recall: $\card{G}=p^{e}m$. Let $s$ be the number of $p$-Sylow subgroups. Then $s\mid{}m$ and $s\equiv{}1\pmod{p}$.\n
}

\par\noindent
Remark: Let $H$ be a $p$-Sylow subgroup of $G$. Then $H$ is a normal subgroup if and only if $s=1$. Can often prove $s=1$ using the third Sylow theorem.\n
Now, we move on to prove the theorems...\n

\newpage

\section*{First Sylow Theorem Proof}

\par\noindent
Let $\card{G}=p^{e}m$. let $\mathcal{U}$ be the set of all subsets of $G$ of size $p^{e}$.\n
$G$ acts on $\mathcal{U}$: given $g\in{}G$, $U\in\mathcal{U}$, $g\cdot{}U=gU=\set{gh\mid{}h\in{}U}$.\n

\prop{
	$\#\mathcal{U}=\binom{n}{p^{e}}=\frac{n(n-1)(n-2)\cdots(n-p^{e}+1)}{p^{e}(p^{e}-1)\cdots1}$. This is a general fact: the number of $k$-element subsets of a set of size $n$ is $\binom{n}{k}$.\n
}

\prop{
	$p\ndiv\#\mathcal{U}$.\n
	Proof: $\#\mathcal{U}=\frac{n(n-1)\cdots(n-p^{e}+1)}{p^{e}(p^{e}-1)\cdots1}$. If $0\le{}k<p^{e}$, then the power of $p$ dividing $n-k$ is the same as $p^{e}-k$, so the $p$'s in the numerator and denominator cancel in each pair.\n
	$\#\mathcal{U}=\#O_{1}+\cdots+\#O_{r}$, where the $O_{i}$'s are the orbits of $G$. Since $p\ndiv\#\mathcal{U}$, we must have $p\ndiv\#O_{i}$ for some $i$. Say $O_{i}$ is the orbit of $U\in\mathcal{U}$. Let $H=\stab(u)$. $\forall{}h\in{}H,x\in{}U$, we have $hx\in{}U$, so $U$ contains the coset $Hx$. Thus, $U$ is a union of cosets, so $\#H\mid\#U=p^{e}$. Thus, $\#H=p^{k}$ for some $0\le{}k\le{}e$. By the counting theorem, $\#O_{i}\cdot\#\stab(u)=\#G=p^{e}m$. $\#\stab(u)=\#H=p^{k}$, so $\#O_{i}=p^{e-k}\cdot{}m$. Since $p\ndiv\#O_{u}$, $e=k$, so $H$ is a $p$-Sylow.\proven
}

\section*{Second Sylow Theorem Proof}

\par\noindent
Let $G$ be a group, $H\subseteq{}G$ a $p$-Sylow, and $K\subseteq{}G$ some other subgroup. We want to show there's a conjugate of $H$, $H'$, \st{} $H'\cap{}K$ is a $p$-Sylow of $K$.\n

\par\noindent
Consider the action of $G$ on $G/H$. Recall that the stabilizers for this action are conjugates of $H$.\n
$\#G/H=m$, so $p\ndiv\#G/H$. Thus, there exists an orbit of $K$ on $G/H$ of cardinality not divisible by $p$. Say its the orbit of $gH$. The stabilizer of $gH$ in $G$ is $H'=gHg\inv$, so the stabilizer in $K$ is $H'\cap{}K$.\n

\par\noindent
$H'\cap{}K\subseteq{}H'$, so $H'\cap{}K$ is a $p$-group. By the counting formula (for the action of $K$ on $G/H$),\n
\[
	\#O_{gH}\cdot\underbrace{\#\stab(gH)}_{H'\cap{}K}=\#K\quad\Rightarrow\quad\#O_{gH}=\frac{\#K}{\#H'\cap{}K}\quad\Rightarrow\quad{}p\ndiv\frac{\#K}{\#H'\cap{}K}\quad\Rightarrow\quad{}H'\cap{}K\ptxt{ is a $p$-Sylow}
\]
\proven

\section*{Third Sylow Theorem Proof}

\par\noindent
Assume $\card{G}=p^{e}m$. Let $s$ be the number of $p$-Sylow subgroups of $G$. We need to show (1) $s\mid{}m$ and (2) $s\equiv{}1\pmod{p}$.\n

\begin{enumerate}[label=(\arabic*)]
	\item Let $\mathcal{H}=\set{H_{1},\ldots,H_{s}}$ be the set of all $p$-Sylows. $G$ acts on $\mathcal{H}$ by conjugation. This action is transitive by a corollary of the second theorem. Let $H=H_{1}$. What is $\stab(H)$? It's $\set{g\mid{}gHg\inv=H}=N$, called the \textbf{normalizer} of $H$. It's clear that $H\subset{}N$. By the counting formula,
	\[
		\#\underbrace{O_{H}}_{\mathcal{H}}\cdot\#\underbrace{\stab_{H}}_{N}=\#G
	\]
	So $s=\frac{\#G}{\#N}$, and $\#G=p^{e}m$, and $p^{e}=\#H\mid\#N$, so we conclude $s\mid{}m$.\up\proven
	\item Think about $H$ acting on $\mathcal{H}$ by conjugation. $H$ fixes $H=H_{1}$. We claim this is the only fixed point.\n
	Proof: Suppose $H$ fixes $H_{i}$. This implies $H\cdot{}H_{i}$ is a subgroup, because for $ab\in{}HH_{i}$, $a'b'\in{}HH_{i}$, $(ab)(a'b')=\underbrace{(aa')}_{\in{}H}\underbrace{(a')\inv{}ba'b'}_{\in{}H_{i}}$.\n
	Exercise: $HH_{i}$ is a $p$-group. Since $H\cdot{}H_{i}$ contains $H$, we must have $H=H\cdot{}H_{i}$.\n
	So $H=H_{i}$, thus $i=1$.\n
	Now, use the class equation for $H\acts\mathcal{H}$.
	\begin{itemize}
		\item $O_{H_{1}}$ has size $1$.
		\item Every other orbit has size divisible by $1$.
	\end{itemize}
	Thus, $s=1\pmod{p}$.\up\proven
\end{enumerate}

\subsection*{Groups of Order $15$, $21$, $12$}

\prop{
	Every group of order $15$ is cyclic.\n
	Proof: Let $H$ be a $3$-Sylow, $K$ be a $5$-Sylow. $H\cong\Z/3\Z$, $K\cong\Z/5\Z$.\n
	Let $s$ be the number of $3$-Sylow subgroups. By the third Sylow theorem, $s\mid{}5$, so $s=1$ or $s=5$, and $s\equiv{}1\pmod{3}$. We must have $s=1$. Thus, $H$ is the unique $3$-Sylow, so $H$ is normal.\n
	Let $s'$ be the number of $5$-Sylow subgroups. $s'\mid{}3$, so $s'=1$ or $s'=3$. $s'\equiv{}1\pmod{5}$, so $s'=1$. Thus $K$ is normal.\n
	\n
	We claim that $G=HK$. Proof: we know $HK$ is a subgroup, because $H$ is normal, and that it contains $H$ and $K$. Thus, $3,5\mid\#HK$, and $\#HK\mid{}\#G=15$. So $\#HK=15$. Clearly, $H\cap{}K=\set{1}$, because $\#H\cap{}K\mid\gcd(\#H,\#K)=1$. So $G\equiv{}H\times{}K\cong\Z/15\Z$.\proven
}

\prop{
	There are $2$ groups of order $21$ up to isomorphism.\n
	Proof: Let $G$ be a group of order $21=3\cdot{}7$. Let $s$ be the number of $7$-Sylows. By the third theorem, $s\mid{}3$, so $s=1$ or $3$. Because $s\equiv{}1\pmod{7}$, $s=1$. Let $H$ be the unique $7$-Sylow (note that it's normal). Let $K$ be a $3$-Sylow. Just as in the previous proof, we know $G=H\cdot{}K$ and $H\cap{}K=\set{1}$. So $G\cong{}H\rtimes{}K$.\n
	The structure of the semi-direct product is determined by the action of $K$ on $H$.
	\[
		\Z/3\Z\cong{}K\to\Aut(H)=\F_{7}^{\times}\cong\Z/6\Z\overset{\ptxt{CRT}}{\cong}\Z/3\Z\times\Z/2\Z
	\]
	Thus, we have two groups of order $21$ (up to isomorphism):
	\begin{itemize}[leftmargin=4\parindent]
		\item $\Z/21\Z$
		\item $\Z/7\Z\rtimes_{\varphi}\Z/3\Z$ where $\varphi:\Z/3\Z\to\Aut(\Z/7\Z)$ is nontrivial.
	\end{itemize}\up\n
	\proven
}

\prop{
	There are $5$ groups of order $12$ up to isomorphism. These groups are
	\begin{itemize}[leftmargin=4\parindent]
		\item $\Z/12\Z=\Z/3\Z\times\Z/4\Z$
		\item $\Z/6\Z\times\Z/2\Z=\Z/3\Z\times\Z/2\Z\times\Z/2\Z$
		\item $A_{4}$ (the alternating group)
		\item $D_{6}$ (the dihedral group)
		\item One more
	\end{itemize}\up\n
	Proof: Let $H$ be a $2$-Sylow, so $H=\Z/4\Z$ or $H=\Z/2\Z\times\Z/2\Z$.\n
	Let $K$ be a $3$-Sylow, so $K\cong\Z/3\Z$.\n
	The key claim we will make is that at least one of $H$ or $K$ is normal.\n
	Proof: Let $s$ be the number of $3$-Sylows. $s\mid{}4$, and $s\equiv{}1\pmod{3}$, so $s=1$ or $s=4$. If $s=1$, then $K$ is normal. \textit{I'm missing the remainder of the proof from my notes.}\n
	Thus, we have $G=HK$, with $H\cap{}K=\set{1}$. So $G\cong{}H\rtimes{}K$ or $G\cong{}K\rtimes{}H$.\proven
}

\end{document}